\documentclass[11pt]{article}
\usepackage[T1]{fontenc}
\usepackage{lmodern,amsmath,amssymb,amsthm,booktabs,microtype}
\usepackage[hmargin=0.96in,vmargin=0.84in]{geometry}
\setlength{\emergencystretch}{1.5em}
\usepackage[colorlinks=true,linkcolor=blue,urlcolor=blue,citecolor=blue]{hyperref}
\newtheorem{theorem}{Theorem}
\newtheorem{lemma}[theorem]{Lemma}
\newcommand{\R}{\mathbb R}
\newcommand{\Z}{\mathbb Z}
\newcommand{\Lag}{\mathcal L}
\title{Nonreal zeros of Ehrhart numerators for\\three-row transportation polytopes}
\author{}
\date{November 26, 2026}
\begin{document}
\begin{center}
{\LARGE Nonreal zeros of Ehrhart numerators for\\[2pt]three-row transportation polytopes\par}
\medskip November 26, 2026
\end{center}
\begin{abstract}
The transportation polytope with row sums $(12,1,1)$ and column sums $(2,1^{12})$ has an Ehrhart $h^*$-polynomial with a nonreal conjugate pair. This disproves a conjecture of Alexandersson. More generally, for rows $(m,1,1)$ and columns $(2,1^m)$, we prove that the numerator has only real zeros exactly when $2\le m\le11$, among $m\ge2$. Its number of nonreal zeros tends to infinity with $m$. A rational certificate proves the single counterexample independently of the classification.
\end{abstract}

\section{The conjecture and an exact counterexample}
For positive integer vectors $r,c$ of equal sum, $T(r,c)$ is the set of nonnegative real matrices with row sums $r$ and column sums $c$. It is a lattice polytope: the constraint matrix becomes a bipartite incidence matrix upon negating the column-sum equations, and is totally unimodular. Writing $d=\dim T$ and $L_T(t)=|tT\cap\Z^{r_0c_0}|$, where $r_0,c_0$ count the rows and columns, its Ehrhart numerator $h^*_T$ is determined by
\[
  \sum_{t\ge0}L_T(t)z^t=\frac{h^*_T(z)}{(1-z)^{d+1}}.
\]
Alexandersson conjectured real-rootedness whenever each margin vector contains an entry at least two \cite{symcat}.

\begin{theorem}\label{thm:counter}
Set $T_m=T((m,1,1),(2,1^m))$ for $m\ge2$, where $1^m$ denotes $m$ entries equal to one. Then
\begin{equation}\label{eq:h}
 H_m(z):=h^*_{T_m}(z)=B_m(z)-mz,\qquad B_m(z)=\sum_{j=0}^m\binom mj^2z^j.
\end{equation}
In particular, $H_{12}$ has a nonreal zero.
\end{theorem}
\begin{proof}
The affine constraints have rank $3+(m+1)-1=m+3$. A strictly positive feasible table is the outer product of the two margins divided by their common total. Thus $\dim T_m=2m$.

In an integral point of $tT_m$, the two rows of sum $t$ determine the first row. Their unrestricted choices are two weak compositions of $t$ into $m+1$ parts. In the column of capacity $2t$ every pair fits. A unit-margin column fails precisely when the two entries there sum to at least $t+1$. Two such failures are impossible because the rows together have sum $2t$. For a specified failed column, its $2m$ off-column entries have total at most $t-1$. Conversely, each such choice gives a unique failed column, since both remaining row sums are nonnegative and their sum is at least $t+1$. Stars and bars now gives
\begin{equation}\label{eq:count}
 L_{T_m}(t)=\binom{t+m}{m}^2-m\binom{t+2m-1}{2m}.
\end{equation}
The second binomial coefficient is zero at $t=0$.

For the first term, differentiate $z^m/(1-z)^{m+1}$ $m$ times and divide by $m!$. Its resulting coefficient of $z^t$ is $\binom{t+m}{m}^2$. Leibniz's rule gives numerator $\sum_{k=0}^m\binom mk\binom{m+k}m z^k(1-z)^{m-k}$ over $(1-z)^{2m+1}$. The coefficient of $z^j$ in this numerator equals $\binom mj$ times the $j$th forward difference of $k\mapsto\binom{m+k}m$ at zero, hence equals $\binom mj^2$ by Pascal's identity. The second term in \eqref{eq:count} has generating function $mz/(1-z)^{2m+1}$ before subtraction. This proves \eqref{eq:h}.

For any real polynomial $H$ with only real zeros $a_i$, repeated with multiplicity, logarithmic differentiation yields
\[
  \Lag_H(x):=H'(x)^2-H(x)H''(x)=H(x)^2\sum_i\frac{1}{(x-a_i)^2}\ge0
\]
away from its zeros and, by continuity, everywhere. At $x=-1/5$, formula \eqref{eq:h} gives
\[
  5^{12}H_{12}(x)=41780056,\quad
  5^{11}H'_{12}(x)=-357481092,\quad
  5^{10}H''_{12}(x)=12374409432.
\]
Consequently $5^{22}\Lag_{H_{12}}(-1/5)=-389210787898375728<0$. Both prescribed margin vectors of $T_{12}$ contain an entry at least two.
\end{proof}

\section{A sharp transition in the family}
\begin{theorem}\label{thm:all}
For integers $m\ge2$, $H_m$ has only real zeros if and only if $2\le m\le11$. For $m\ge256$ it has at least $\lfloor\sqrt m/\pi\rfloor-4$ nonreal zeros counted with multiplicity.
\end{theorem}

\begin{lemma}\label{lem:osc}
For $m\ge256$ let $I_m=[1/(2m),2/m]$ and $q=\lfloor\sqrt m/\pi\rfloor$. Then $H_m(-y)>0$ on $I_m$, while $B_m(-y)$ has at least $q$ distinct zeros in its interior.
\end{lemma}
\begin{proof}
Write $y=\tan^2(\theta/2)$. The finite binomial formula for the Legendre polynomial $P_m$, normalized by $P_m(1)=1$, gives \cite[Eq.~18.5.8]{dlmf}
\[
 (1+y)^{-m}B_m(-y)=P_m\!\left(\frac{1-y}{1+y}\right)=P_m(\cos\theta).
\]
The Bernstein bound \cite[Eq.~18.14.3, $\alpha=\beta=0$]{dlmf} gives
\[
 |P_m(\cos\theta)|\le\sqrt{\frac{2}{\pi(m+1/2)\sin\theta}}.
\]
On $I_m$ one has $\sin\theta=2\sqrt y/(1+y)\ge \sqrt{2/m}/(1+2/m)\ge1/\sqrt m$. Thus $|P_m(\cos\theta)|<m^{-1/4}\le1/4$. With $s=my\in[1/2,2]$, the correction satisfies
\[
   \frac{my}{(1+y)^m}\ \ge\ se^{-s}\ \ge\ \min\{e^{-1/2}/2,2e^{-2}\}\ >\ 1/4.
\]
This proves positivity.

Set $M=m+1/2$. The Legendre differential equation, after the substitution $u(\theta)=\sqrt{\sin\theta}\,P_m(\cos\theta)$, becomes
\[
 u''+\left(M^2+\frac{1}{4\sin^2\theta}\right)u=0.
\]
Every closed interval $[a,b]\subset(0,\pi)$ of length $\pi/M$ contains an interior zero of $u$. Indeed, if $u>0$ on its interior (multiply by $-1$ if needed), put $v(\theta)=\sin M(\theta-a)$. Then $(uv'-u'v)'=uv/(4\sin^2\theta)>0$ there. The endpoint difference is $-M(u(a)+u(b))\le0$, a contradiction.

The $\theta$-interval corresponding to $I_m$ has length
\[
 2\!\int_{1/\sqrt{2m}}^{\sqrt{2/m}}\!\frac{dt}{1+t^2}
 \ \ge\ \frac{\sqrt2}{(1+2/m)\sqrt m}\ >\ \frac1{\sqrt m}\ >\ \frac{q\pi}{M}.
\]
It contains $q$ intervals of length $\pi/M$ with disjoint interiors. Their interior zeros are distinct and yield $q$ zeros of $B_m(-y)$.
\end{proof}

\begin{proof}[Proof of Theorem~\ref{thm:all}]
First take $m\ge256$. Since $H_m(-y)-B_m(-y)=my$, the two functions have the same second derivative in $y$. Lemma~\ref{lem:osc} and two applications of Rolle's theorem give at least $q-2$ zeros of this derivative inside $I_m$. Let $N_L,N_R$ be the real zero counts, with multiplicity, of $H_m(-y)$ to the left and right of $I_m$. The same derivative has at least $N_L-2$ zeros on the left and $N_R-2$ on the right, by Rolle's theorem with multiplicities. Its degree is $m-2$, so
\[
 m-2\ge(q-2)+(N_L-2)+(N_R-2).
\]
There are no zeros of $H_m(-y)$ on $I_m$ by Lemma~\ref{lem:osc}. Hence the number of nonreal zeros is $m-N_L-N_R\ge q-4$, which is positive already at $m=256$.

For the finite ranges the following sign certificates use only integer arithmetic. For a polynomial $f$ of degree $d$ put
\[
 E(f,p,s)=\sum_{k=0}^d [z^k]f\;p^k s^{d-k}\qquad(s>0).
\]
For $m=2$, $H_m=(1+z)^2$. For $3\le m\le11$, evaluate $E(H_m,-j,256-j)$ for $0\le j<256$, discard zeros, and append the sign $(-1)^m$ at negative infinity. The numbers of sign changes, in order of increasing $m$, are $3,4,5,6,7,8,9,10,11$. The arguments are strictly decreasing, so these signs bracket $m$ distinct negative zeros in each case.

For $12\le m\le255$, use $s=6m$ and $1\le j\le18$. Put
\[
 C_{m,j}=E(H'_m,-j,s)^2-E(H_m,-j,s)E(H''_m,-j,s).
\]
For all 244 values of $m$, $\min_j C_{m,j}<0$. The following finite loop and accompanying implementation reproduce this last certificate.
\begin{verbatim}
for m = 12,...,255:
    h[k] = binomial(m,k)^2 for k = 0,...,m
    h[1] = h[1] - m
    f = derivative(h); g = derivative(f)
    require some j in {1,...,18} satisfies
        E(f,-j,6*m)^2 < E(h,-j,6*m) * E(g,-j,6*m)
\end{verbatim}
Since $C_{m,j}=(6m)^{2m-2}\Lag_{H_m}(-j/(6m))$, each such certificate proves that $H_m$ has a nonreal zero. This completes the classification.
\end{proof}

\paragraph{Verification.} The file \texttt{verify.py} performs every finite computation in the proof. For $m=12$ it also recovers every numerator coefficient independently: it fills the thirteen columns successively, retaining the cumulative totals in the two unit rows, and multiplies the resulting counting series by $(1-z)^{25}$ through degree 24. The two methods agree in every coefficient.

\begin{thebibliography}{9}
\bibitem{symcat} Per Alexandersson. \emph{Real-rooted polynomials: open problems}. The Symmetric Functions Catalog, entry \texttt{conjTransportationHStar}, consulted November 26, 2026. \url{https://www.symmetricfunctions.com/realRootedOpenProblems.htm#conjTransportationHStar}.
\bibitem{dlmf} F. W. J. Olver et al., editors. \emph{NIST Digital Library of Mathematical Functions}, Chapter 18. \url{https://dlmf.nist.gov/18.5.E8} and \url{https://dlmf.nist.gov/18.14.E3}.
\end{thebibliography}
\end{document}
